\subsection{LoRaWAN Scalability and the ``SF Monoculture''}

LoRaWAN is the de facto standard for long-range, low-power IoT
connectivity~\cite{augustin2016lora,bor2016lora}. Augustin et al.~\cite{augustin2016lora}
provide a foundational analysis of the LoRa physical layer, characterising the chirp
spread-spectrum modulation, the six orthogonal spreading factors (\(\text{SF} \in \{7,
\ldots, 12\}\)), and the trade-off between data rate and link budget that governs SF
selection. Bor et al.~\cite{bor2016lora} demonstrate through large-scale simulation that a
single LoRaWAN gateway can theoretically serve thousands of devices, but that this capacity
is highly sensitive to the SF distribution: if too many devices converge on the same
high-SF setting, per-device throughput collapses due to intra-SF collisions and the
resulting long Time-on-Air. These two foundational works establish the protocol constraints
that every LoRaWAN system designer---and by extension every UAV mission planner---must
respect.

\subsubsection{The SF Monoculture and ADR Failure}

A critical failure mode in mobile LoRaWAN networks is the congestion collapse known as the
``Spreading Factor (SF) Monoculture'' or ``SF12 Well.'' Zorbas et
al.~\cite{zorbas2025revisiting} provide a rigorous mathematical analysis of this phenomenon,
demonstrating that the standard LoRaWAN Adaptive Data Rate (ADR) mechanism---which adjusts
each device's SF based on a smoothed RSSI history---fails systematically in dense
deployments. Because ADR reacts to past link quality rather than current conditions, devices
in a rapidly changing environment (such as one serviced by a moving UAV) accumulate
outdated RSSI estimates. This latency in the ADR control loop causes devices to converge on
high SF settings long after the UAV has moved away, eventually driving the entire network
toward SF12---the slowest, longest Time-on-Air setting. The consequence is a self-reinforcing
congestion collapse: high SF increases collision probability, reducing successful
transmissions, which in turn prevents devices from receiving ADR downgrade commands. Zorbas
et al.\ conclude that ``reinforcement learning models could be employed to adapt over time
to evolving network conditions,'' a direct motivation for the DRL approach in this
dissertation.

Reynders et al.~\cite{reynders2016range} provide complementary evidence: collision
probability grows super-linearly beyond a critical load threshold, and frequency-adaptive
rate management is \emph{necessary} to sustain performance, motivating the ADR modelling
choices in this work.

\subsubsection{The Capture Effect as a Partial Remedy}

An important physical mitigation mechanism in congested LoRaWAN networks is the LoRa
Capture Effect. Croce et al.~\cite{croce2020lora} are the first to rigorously
characterise this phenomenon through controlled software-defined radio (SDR) experiments.
They demonstrate experimentally that collisions between two LoRa signals on the \emph{same}
spreading factor can be resolved if the stronger signal exceeds the weaker by a power ratio
of approximately 6\,dB---the ``capture threshold.'' Below this threshold, both packets are
lost; above it, the stronger packet is decoded successfully. This finding overturns a
common simplifying assumption in prior LoRaWAN simulators (that any intra-SF collision
results in mutual destruction) and has significant implications for capacity analysis.

Croce et al.\ further show that, contrary to intuition, allocating high SFs to far-away
devices (to compensate for path loss) can actually \emph{reduce} cell capacity in congested
deployments, because high-SF packets occupy the channel for longer, increasing the
probability that a second device transmits during the same Time-on-Air window. This creates
a direct tension between coverage (which demands high SFs at range) and capacity (which
demands low SFs for throughput). A mobile UAV that can close the physical distance to
sensors---reducing path loss and enabling SF downgrade via ADR---can resolve this tension
dynamically, something no static gateway can achieve.

This work directly incorporates the 6\,dB threshold into the collision resolution logic,
rewarding the agent for reaching spatial positions that create the power differential
required for capture---a design distinction from all protocol-blind trajectory planners.

